\documentclass[11pt]{article}
\usepackage[margin=1in]{geometry}
\usepackage{amsmath,amssymb,amsthm,mathtools}
\usepackage{graphicx}
\usepackage{microtype}
\usepackage[hidelinks]{hyperref}
\usepackage{enumitem}

\newtheorem{theorem}{Theorem}
\newtheorem{lemma}{Lemma}
\newtheorem{corollary}{Corollary}
\newtheorem{remark}{Remark}
\newcommand{\M}{\mathcal M}
\newcommand{\A}{\mathcal A}
\newcommand{\B}{\mathcal B}
\newcommand{\N}{\mathcal N}
\newcommand{\Z}{\mathcal Z}
\newcommand{\D}{\mathcal D}
\newcommand{\E}{\mathbb E}
\newcommand{\C}{\mathbb C}
\makeatletter
\renewcommand{\@seccntformat}[1]{\csname the#1\endcsname\hspace{0.65em}}
\makeatother

\title{Central branching is the only maximality obstruction\\
for zero-diagonal MASAs in a type $\mathrm{II}_1$ factor}
\author{Automated mathematical discovery packet}
\date{August 2026}

\begin{document}
\maketitle

\begin{center}
\fbox{\parbox{0.91\textwidth}{\textbf{Status.}
Candidate partial result, likely valid, subject to human review. The full
MASA question remains open.}}
\end{center}

\section{Source problem}

The source is K.~Dykema and P.~Skoufranis, \emph{Numerical Ranges in
$\mathrm{II}_1$ Factors}, arXiv:1503.05766 \cite{DS}. On physical PDF page 19,
immediately before Lemma 6.1, the authors write:

\begin{quote}
``Unfortunately, we do not obtain true `diagonals' as we do not know if one
can guarantee $\A$ in the following technical lemma (whose proof is a
generalization of a matricial result) is a MASA.''
\end{quote}

The lemma then says that if $(\M,\tau)$ is a type $\mathrm{II}_1$ factor and
$\tau(T)=0$, there is a diffuse abelian von Neumann subalgebra $\A\subset\M$
such that $E_{\A}(T)=0$. Thus the question is:

\begin{quote}
\textbf{Can $\A$ always be chosen maximal abelian?}
\end{quote}

\begin{figure}[ht]
\centering
\includegraphics[width=\textwidth]{figures/open_problem_crop.png}
\caption{Exact source passage and Lemma 6.1, rendered from the stored arXiv
source (physical PDF page 19).}
\end{figure}

\section{Definitions}

All conditional expectations are the unique trace-preserving normal
expectations. Call an abelian von Neumann algebra $\A\subset\M$
\emph{zero-diagonal for $T$} if $E_{\A}(T)=0$. This is equivalent to
\[
   \tau(Ta)=0 \qquad (a\in\A).
\]
For such $\A$, set
\[
 \N=\A'\cap\M,\qquad \Z=Z(\N),\qquad g=E_{\Z}(T).
\]
Since every element of $\N$ commutes with $\A$, one has $\A\subset\Z$.

We use the following concrete form of relative atomlessness. The inclusion
$\A\subset\Z$ is \emph{relatively atomless} if, after identifying the two
abelian algebras with nested sub-$\sigma$-algebras of a standard probability
space, the larger measure algebra has no atom relative to the smaller one.
The inclusion
$\A\subset\Z$ has the \emph{conditional halving property for $g$} if there is
a projection $q\in\Z$ such that
\[
 E_{\A}(q)=\tfrac12 1,
 \qquad E_{\A}(gq)=0.
\]
The conditional Lyapunov theorem \cite{DE} implies this property whenever the
abelian extension $\A\subset\Z$ is relatively atomless.

\section{Proof intuition}

First maximize, by inclusion, an abelian algebra on which $T$ has zero
diagonal. Such a maximal algebra cannot have an atom: the source lemma can be
inserted into the corresponding $\mathrm{II}_1$ corner and gives a strict
zero-diagonal enlargement.

The only place where maximality might still fall short of being a MASA is the
relative commutant $\N=\A'\cap\M$. If the central part of $T$ in $\N$ vanishes,
then every factor fiber of $\N$ contains a zero-diagonal abelian algebra. A
measurable choice of these fiber algebras assembles to a zero-diagonal
extension containing the center, contradicting maximality. Hence a non-MASA
maximal algebra must carry a nonzero central remnant $g$. Even that remnant can
be split away when the central extension is relatively atomless. The sole
remaining configuration is therefore nonzero, conditionally atomic central
branching.

\section{The center-valued zero-diagonal lemma}

\begin{lemma}[center-valued zero diagonal]\label{lem:center}
Let $(\N,\tau)$ be a finite von Neumann algebra with separable predual, let
$\Z=Z(\N)$, and let $S\in\N$ satisfy $E_{\Z}(S)=0$. Then there is an abelian
von Neumann algebra $\D$ such that
\[
 \Z\subset\D\subset\N,
 \qquad E_{\D}(S)=0.
\]
If $\N$ is nonabelian on a nonzero central summand, $\D$ can be chosen to
strictly contain $\Z$ on that summand.
\end{lemma}

\begin{proof}
Decompose $\N$ over its center as a standard measurable field of finite
factors,
\[
  \N=\int_X^{\oplus}\N_x\,d\mu(x),
  \qquad S=\int_X^{\oplus}S_x\,d\mu(x).
\]
The hypothesis says that the normalized scalar trace $\tau_x(S_x)$ is zero
for almost every $x$.

If $\N_x=M_n(\C)$, choose a unit vector $e_1$ with
$\langle S_xe_1,e_1\rangle=0$. Such a vector exists because the numerical
range is convex, contains the spectrum, and therefore contains the average of
the eigenvalues, namely zero. Compressing to $e_1^\perp$ leaves trace zero;
induction gives an orthonormal basis in which every diagonal entry is zero.
The corresponding diagonal algebra $\D_x$ satisfies $E_{\D_x}(S_x)=0$.
For a type $\mathrm{II}_1$ fiber, Lemma 6.1 of \cite{DS} supplies a diffuse
abelian $\D_x$ with the same property. A scalar fiber has $S_x=0$, and we put
$\D_x=\C$ there.

It remains to make the choices measurable. On the unit balls of a separable
measurable field, the projection relation, normalized trace, compression, and
the equations $\tau_x(S_xp)=0$ are Borel conditions and are closed in the
$2$-norm on bounded sets. The nonempty solution fields therefore admit
measurable selectors by the Kuratowski--Ryll-Nardzewski theorem
\cite{KRN}. Applying this recursively to the finite-dimensional basis choices
and to the dyadic projection tree in the proof of Lemma 6.1 produces a
measurable field $x\mapsto\D_x$. Direct-integral theory \cite{Tak} then gives
\[
 \D=\int_X^{\oplus}\D_x\,d\mu(x).
\]
It contains the diagonal center $L^\infty(X)$, and the fiberwise identities
give $E_{\D}(S)=0$. On every nonabelian factor fiber the construction contains
a nonscalar projection, so measurable selection on any positive-measure
nonabelian stratum makes $\D$ strictly larger than $\Z$ there.
\end{proof}

\section{Main partial theorem}

\begin{theorem}[central-obstruction reduction]\label{thm:main}
Let $(\M,\tau)$ be a type $\mathrm{II}_1$ factor with separable predual and
let $T\in\M$ satisfy $\tau(T)=0$.

\begin{enumerate}[label=(\roman*)]
\item There is an abelian zero-diagonal algebra $\A$ maximal by inclusion,
and every such maximal $\A$ is diffuse.
\item For a maximal $\A$, put $\N=\A'\cap\M$ and $\Z=Z(\N)$. If
$E_{\Z}(T)=0$, then $\A$ is a MASA.
\item In particular, $Z(\A'\cap\M)=\A$ implies that $\A$ is a MASA.
\item If $\A$ is not a MASA, then $\A\subset\Z$ fails the conditional halving
property for $g=E_{\Z}(T)$. Consequently, a maximal $\A$ whose strict central
extension is relatively atomless must be a MASA.
\end{enumerate}

Thus, if a maximal zero-diagonal $\A$ is not a MASA, then necessarily
\[
 \A\subsetneq Z(\A'\cap\M),\qquad E_{Z(\A'\cap\M)}(T)\ne0,
\]
and the central extension fails relative atomlessness.
\end{theorem}

\begin{proof}
Lemma 6.1 of \cite{DS} makes the family of zero-diagonal abelian algebras
nonempty. Let $(\A_i)$ be a chain in this family and put
$\A=(\bigcup_i\A_i)''$. Every $a\in\A$ is an $L^2$ limit, uniformly bounded
after a standard truncation, of elements from the union. Hence
$\tau(Ta)=0$, so $E_{\A}(T)=0$. Zorn's lemma gives a maximal $\A$.

Suppose that $p$ is an atom of $\A$. Since $E_{\A}(T)=0$,
$\tau(pTp)=\tau(Tp)=0$. The corner $p\M p$ is a type $\mathrm{II}_1$ factor,
so Lemma 6.1 of \cite{DS} gives a diffuse abelian algebra
$\B_p\subset p\M p$ with $E_{\B_p}(pTp)=0$. Then
\[
   (1-p)\A\oplus\B_p
\]
strictly contains $\A$ and still pairs to zero with $T$, contradicting
maximality. Thus $\A$ is diffuse. This proves (i).

Put $\N=\A'\cap\M$, $\Z=Z(\N)$, and $S=E_{\N}(T)$. Conditional-expectation
composition gives
\[
 E_{\Z}(S)=E_{\Z}(T).
\]
If this is zero, Lemma~\ref{lem:center} gives an abelian
$\D$ with $\Z\subset\D\subset\N$ and $E_{\D}(S)=0$. Since
$E_{\D}(T)=E_{\D}(S)$, maximality forces $\D=\A$. But $\A\subset\Z\subset\D$,
so $\Z=\A$ as well. If $\N$ were larger than $\A$, either it would be
abelian, in which case $\N=\Z$ is already a strict zero-diagonal extension,
or it would be nonabelian on a nonzero central summand, where
Lemma~\ref{lem:center} makes $\D$ strict. Both contradict maximality. Hence
$\N=\A$, exactly the MASA condition. This proves (ii), and (iii) follows from
$E_{\Z}(T)=E_{\A}(T)=0$.

Finally suppose the conditional halving property holds. Let
$\B=\A\vee\{q\}\subset\Z$. Since $E_{\A}(g)=E_{\A}(T)=0$, the two central
atoms satisfy
\[
 E_{\A}(gq)=0,
 \qquad E_{\A}(g(1-q))=0.
\]
Because $E_{\A}(q)=E_{\A}(1-q)=1/2$, the explicit two-atom conditional
expectation gives $E_{\B}(g)=0$. Also
$E_{\B}(T)=E_{\B}(E_{\Z}(T))=E_{\B}(g)=0$. The equality
$E_{\A}(q)=1/2$ forces $q\notin\A$, so $\B$ is a strict zero-diagonal
extension. This contradicts maximality and proves the first assertion of
(iv). Indeed, for a relatively atomless extension, apply the conditional
Lyapunov theorem \cite{DE} to the two-coordinate conditional vector measure
\[
 q\longmapsto \bigl(E_{\A}(q),E_{\A}(gq)\bigr).
\]
Its range contains half of the value at $1$, namely
$\bigl(\tfrac12 1,\tfrac12 E_{\A}(g)\bigr)=(\tfrac12 1,0)$.
Thus it supplies the required projection $q$ and proves the consequence.
\end{proof}

\section{Why the central obstruction cannot simply be deleted}

Consider
\[
 \N=M_2(\C)\oplus\C,
 \quad \tau=\tfrac12\operatorname{tr}_2\oplus\tfrac12\operatorname{id},
 \quad S=I_2\oplus(-1).
\]
Then $\tau(S)=0$. A projection has the form $p\oplus e$, where
$\operatorname{rank}(p)=r\in\{0,1,2\}$ and $e\in\{0,1\}$. Directly,
\[
 \tau(S(p\oplus e))=\frac r4-\frac e2.
\]
This vanishes only for $(r,e)=(0,0)$ or $(2,1)$, namely the zero and identity
projections. Thus scalar trace zero in a finite algebra does not force a
nontrivial zero-diagonal abelian algebra when the center-valued trace is
nonzero. This is an obstruction model, not a counterexample to the factor
question.

\section{Verification and eight upgrade attempts}

The proof was checked step by step in \texttt{VERIFIER\_REPORT.md}. The finite
obstruction script exhausts all six projection-rank cases and prints
\texttt{VERDICT: PASS}. The eight focused routes were:
\begin{enumerate}
\item Zorn maximalization and atom removal (successful);
\item arbitrary MASA extension (fails because expectation need not persist);
\item localization to the relative commutant (successful reduction);
\item measurable factor-fiber construction under $E_{\Z}(T)=0$ (successful);
\item forcing $Z(\A'\cap\M)=\A$ (no preserving genericity argument found);
\item conditional Lyapunov splitting (successful refinement);
\item direct elimination of the central term (blocked by the finite model);
\item commutator/counterexample routes and a bounded literature audit (no full
resolution found).
\end{enumerate}
The detailed attempt log is stored in the run's \texttt{attempts/} directory.

\section{Literature and novelty check}

The bounded search covered the exact title and arXiv id, the exact phrase
``we do not know if one can guarantee'', the author pair, and combinations of
``trace zero'', ``zero diagonal'', ``MASA'', ``conditional expectation'', and
``type II$_1$ factor'' in arXiv and primary publication pages. Searches also
covered later Schur--Horn and multivariable Schur--Horn papers. No paper
explicitly resolving the source question was found. Marcoux's 2022 survey
still states the broader C-star-algebra version as Problem 3.3 \cite{Mar}.

Novelty confidence is moderate. The ingredients of
Lemma~\ref{lem:center} are standard, but this precise maximality/central
branching reduction was not found in the bounded search.

\section{Limitations and human-review recommendation}

The theorem is partial and assumes separable predual in the measurable
selection step. It neither proves that the central expectation must vanish nor
constructs a factor counterexample when it does not. Human review should focus
on the Effros-Borel assembly in Lemma~\ref{lem:center} and on the precise
conditional Lyapunov hypothesis in part (iv). If those standard measurable
steps are accepted, the rest of the proof is elementary.

\begin{thebibliography}{9}
\bibitem{DS}
K.~Dykema and P.~Skoufranis,
\emph{Numerical Ranges in $\mathrm{II}_1$ Factors},
Proc. Edinburgh Math. Soc. 61 (2018), 31--55;
arXiv:1503.05766.

\bibitem{KRN}
K.~Kuratowski and C.~Ryll-Nardzewski,
\emph{A general theorem on selectors},
Bull. Acad. Polon. Sci. Ser. Sci. Math. Astronom. Phys. 13 (1965), 397--403.

\bibitem{Tak}
M.~Takesaki,
\emph{Theory of Operator Algebras I},
Springer, 1979, direct-integral chapters.

\bibitem{DE}
E.~B.~Dynkin and I.~V.~Evstigneev,
\emph{Regular conditional expectations of correspondences},
Theory Probab. Appl. 21 (1977), 325--338;
doi:10.1137/1121037.

\bibitem{Mar}
L.~W.~Marcoux,
\emph{Five Hilbert Space Problems in Operator Algebras},
Complex Anal. Oper. Theory 16 (2022), article 116.
\end{thebibliography}

\end{document}
